% GENERATED FILE. Edit canonical lessons, then run npm run generate:books.
% canonical-source-hash: fnv1a64:1c1e4794a700db73
% canonical-lessons: ZH-C70-rou, ZH-C70-baozi, ZH-C70-miantiao, ZH-C70-jiaozi, ZH-C70-tang

\chapter{More to Eat}
\label{ch:zh-more-to-eat}
\hlchaptermodality{70}

\begin{chapteropening}
\textbf{By the end of this chapter:} \emph{I can name meat, buns, noodles, dumplings and soup.}

Complete the last lesson of chapter 70: I can name meat, buns, noodles, dumplings and soup.
\end{chapteropening}

\section[ròu]{ròu — meat}

\label{lesson:ZH-C70-rou}



\begin{quote}
Before the new one: say the Mandarin for dinner, then the Mandarin for tasty.
\end{quote}

\subsection*{You'll want to know: ròu}
\emph{ròu} — \textquotedblleft{}meat\textquotedblright{}.

\begin{quote}
\emph{wǒ bù chī ròu} — I don't eat meat.
\end{quote}

Its characters come later; for now it is a word for the ear.

\begin{cousinweb}
\emph{ròu} is inherited from Middle Chinese \emph{nyuwk} (Guangyun, in Baxter's transcription), glossed with the pair 'bone and flesh'.
\end{cousinweb}

\begin{grammarlens}[title={What you've built}]
One more word for everyday answers.
\end{grammarlens}

\subsection*{Your turn}
\begin{itemize}
\raggedright
  \item \emph{Say it:} \emph{ròu}
  \item \emph{Say it:} \emph{ròu}, once more
  \item \emph{Say it:} \emph{wǒ bù chī ròu} — I don't eat meat.
  \item \emph{Recall:} say the Mandarin for dinner, then the Mandarin for tasty, then say \emph{ròu} again
\end{itemize}

\begin{tcolorbox}[breakable,colback=teal!4,colframe=teal!35!black,title={Before you move on}]
What does ròu mean? (\textquotedblleft{}Meat\textquotedblright{}.) Say it once more.
\end{tcolorbox}
\section[bāozi]{bāozi — steamed bun}

\label{lesson:ZH-C70-baozi}



\begin{quote}
Before the new one: say the Mandarin for tasty, then the Mandarin for meat.
\end{quote}

\subsection*{You'll want to know: bāozi}
\emph{bāozi} — \textquotedblleft{}steamed bun\textquotedblright{}.

\begin{quote}
\emph{wǒ chī le sān ge bāozi} — I ate three buns.
\end{quote}

Its characters come later; for now it is a word for the ear.

\begin{cousinweb}
\emph{bāozi} is built from bāo 'to wrap', the same bāo as the word for bag, and the noun ending zi, both inherited from Middle Chinese, \emph{paew} and \emph{tsiX} (Guangyun, in Baxter's transcription): a filling wrapped in dough.
\end{cousinweb}

\begin{grammarlens}[title={What you've built}]
One more word for everyday answers.
\end{grammarlens}

\subsection*{Your turn}
\begin{itemize}
\raggedright
  \item \emph{Say it:} \emph{bāozi}
  \item \emph{Say it:} \emph{bāozi}, once more
  \item \emph{Say it:} \emph{wǒ chī le sān ge bāozi} — I ate three buns.
  \item \emph{Recall:} say the Mandarin for tasty, then the Mandarin for meat, then say \emph{bāozi} again
\end{itemize}

\begin{tcolorbox}[breakable,colback=teal!4,colframe=teal!35!black,title={Before you move on}]
What does bāozi mean? (\textquotedblleft{}Steamed bun\textquotedblright{}.) Say it once more.
\end{tcolorbox}
\section[miàntiáo]{miàntiáo — noodles}

\label{lesson:ZH-C70-miantiao}



\begin{quote}
Before the new one: say the Mandarin for meat, then the Mandarin for steamed bun.
\end{quote}

\subsection*{You'll want to know: miàntiáo}
\emph{miàntiáo} — \textquotedblleft{}noodles\textquotedblright{}.

\begin{quote}
\emph{wǒ xǐhuan chī miàntiáo} — I love noodles.
\end{quote}

Its characters come later; for now it is a word for the ear.

\begin{cousinweb}
\emph{miàntiáo} is built from miàn 'flour' and tiáo 'a twig, a strip', both inherited from Middle Chinese, \emph{menH} and \emph{dew} (Guangyun, in Baxter's transcription): strips of flour.
\end{cousinweb}

\begin{grammarlens}[title={What you've built}]
One more word for everyday answers.
\end{grammarlens}

\subsection*{Your turn}
\begin{itemize}
\raggedright
  \item \emph{Say it:} \emph{miàntiáo}
  \item \emph{Say it:} \emph{miàntiáo}, once more
  \item \emph{Say it:} \emph{wǒ xǐhuan chī miàntiáo} — I love noodles.
  \item \emph{Recall:} say the Mandarin for meat, then the Mandarin for steamed bun, then say \emph{miàntiáo} again
\end{itemize}

\begin{tcolorbox}[breakable,colback=teal!4,colframe=teal!35!black,title={Before you move on}]
What does miàntiáo mean? (\textquotedblleft{}Noodles\textquotedblright{}.) Say it once more.
\end{tcolorbox}
\section[jiǎozi]{jiǎozi — dumpling}

\label{lesson:ZH-C70-jiaozi}



\begin{quote}
Before the new one: say the Mandarin for steamed bun, then the Mandarin for noodles.
\end{quote}

\subsection*{You'll want to know: jiǎozi}
\emph{jiǎozi} — \textquotedblleft{}dumpling\textquotedblright{}.

\begin{quote}
\emph{xīngqītiān wǒmen chī jiǎozi} — We have dumplings on Sunday.
\end{quote}

Its characters come later; for now it is a word for the ear.

\begin{cousinweb}
\emph{jiǎozi} is built from jiǎo and the noun ending zi, which is inherited from Middle Chinese \emph{tsiX} (Guangyun, in Baxter's transcription). The first has no reading in any dictionary we use, so its history is not settled.
\end{cousinweb}

\begin{grammarlens}[title={What you've built}]
One more word for everyday answers.
\end{grammarlens}

\subsection*{Your turn}
\begin{itemize}
\raggedright
  \item \emph{Say it:} \emph{jiǎozi}
  \item \emph{Say it:} \emph{jiǎozi}, once more
  \item \emph{Say it:} \emph{xīngqītiān wǒmen chī jiǎozi} — We have dumplings on Sunday.
  \item \emph{Recall:} say the Mandarin for steamed bun, then the Mandarin for noodles, then say \emph{jiǎozi} again
\end{itemize}

\begin{tcolorbox}[breakable,colback=teal!4,colframe=teal!35!black,title={Before you move on}]
What does jiǎozi mean? (\textquotedblleft{}Dumpling\textquotedblright{}.) Say it once more.
\end{tcolorbox}
\section[tāng]{tāng — soup}

\label{lesson:ZH-C70-tang}



\begin{quote}
Before the new one: say the Mandarin for noodles, then the Mandarin for dumpling.
\end{quote}

\subsection*{You'll want to know: tāng}
\emph{tāng} — \textquotedblleft{}soup\textquotedblright{}.

\begin{quote}
\emph{nǐ hē tāng ma} — Would you like some soup?
\end{quote}

Its characters come later; for now it is a word for the ear.

\begin{cousinweb}
\emph{tāng} is inherited from Middle Chinese \emph{thang} (Guangyun, in Baxter's transcription), glossed 'hot water'. Soup is not given there; from hot water to soup is the likely step, but it is not settled.
\end{cousinweb}

\begin{grammarlens}[title={What you've built}]
One more word for everyday answers.
\end{grammarlens}

\subsection*{Your turn}
\begin{itemize}
\raggedright
  \item \emph{Say it:} \emph{tāng}
  \item \emph{Say it:} \emph{tāng}, once more
  \item \emph{Say it:} \emph{nǐ hē tāng ma} — Would you like some soup?
  \item \emph{Recall:} say the Mandarin for noodles, then the Mandarin for dumpling, then say \emph{tāng} again
\end{itemize}

\begin{tcolorbox}[breakable,colback=teal!4,colframe=teal!35!black,title={Before you move on}]
What does tāng mean? (\textquotedblleft{}Soup\textquotedblright{}.) Say it once more.
\end{tcolorbox}
